\documentclass[10pt]{beamer}
\usepackage[utf8]{inputenc}
\usepackage[T1]{fontenc}
\usepackage[french]{babel}
\usepackage{lmodern}
\usepackage{booktabs}
\usetheme{Madrid}
\usecolortheme{default}

\title{AURA Desktop\newline Presentation technique basee sur le code}
\author{Projet AURA}
\date{11 Mai 2026}

\begin{document}

\begin{frame}
\titlepage
\end{frame}

\begin{frame}{Objectif reel du projet}
\begin{itemize}
\item AURA est une application desktop Python (Tkinter) pour la maintenance industrielle.
\item Elle regroupe en une seule interface: tableaux de bord, analyse NLP, generation d'instructions, planification TRS, orchestration des equipes et historique d'actions.
\item Le systeme est orientee aide a la decision locale (poste Windows), pas une plateforme web distribuee.
\end{itemize}
\end{frame}

\begin{frame}{Architecture logicielle}
\begin{itemize}
\item Point d'entree: \texttt{main.py} lance \texttt{AuraApp}.
\item Interface principale: \texttt{src/gui.py} (onglets metier).
\item Modules metier:
\begin{itemize}
\item \texttt{src/dashboards.py}: radar, gantt, heatmap.
\item \texttt{src/nlp\_engine.py}: extraction manuels + RCA semantique.
\item \texttt{src/prompt\_generator.py}: instructions maintenance + export PDF.
\item \texttt{src/planning.py}: TRS + roadmap 3 ans.
\item \texttt{src/scheduling.py}: affectation techniciens.
\end{itemize}
\item Infrastructure locale:
\begin{itemize}
\item \texttt{src/storage.py}: journal d'evenements SQLite.
\item \texttt{src/app\_logger.py}: logs fichier + console.
\item \texttt{src/app\_config.py}: configuration centralisee.
\end{itemize}
\end{itemize}
\end{frame}

\begin{frame}{Interface utilisateur (ce qui existe)}
\begin{itemize}
\item 6 onglets actifs:
\begin{itemize}
\item Tableau de Bord
\item Analyse NLP
\item Generateur d'Instructions
\item Planification Strategique
\item Orchestration
\item Historique
\end{itemize}
\item Barre de statut en bas de fenetre pour le retour utilisateur.
\item Validation d'entree sur les actions principales (analyse, planification, affectation, etc.).
\end{itemize}
\end{frame}

\begin{frame}{Module Dashboards}
\begin{itemize}
\item Radar: 5 dimensions (Vibrations, Temperature, Pression, Bruit, Lubrification).
\item Gantt: timeline de taches de maintenance (dates + durees).
\item Heatmap: stress/defaillances par composant et jour.
\item Les graphes sont rendus dans Tkinter via Matplotlib (FigureCanvasTkAgg).
\item Simulation possible via generation aleatoire des valeurs radar.
\end{itemize}
\end{frame}

\begin{frame}{Module NLP (exactement selon le code)}
\begin{itemize}
\item Chargement optionnel de:
\begin{itemize}
\item spaCy (modele \texttt{fr\_core\_news\_md})
\item sentence-transformers (\texttt{paraphrase-multilingual-MiniLM-L12-v2})
\end{itemize}
\item Extraction de texte depuis PDF, DOCX, TXT.
\item Extraction d'informations:
\begin{itemize}
\item composants critiques
\item frequences de lubrification
\item seuils d'alerte
\end{itemize}
\item RCA: similarite cosinus entre rapport utilisateur et base de connaissances.
\item Seuil de decision RCA: confiance > 0.4 pour retenir la cause la plus proche.
\end{itemize}
\end{frame}

\begin{frame}{Generateur d'instructions et export}
\begin{itemize}
\item Le generateur n'est pas un LLM distant dans la version actuelle.
\item Mecanisme actuel: recherche par mots-cles (roulement, moteur, pompe) dans un dictionnaire interne.
\item Reponse enrichie par recommendation LOTO.
\item Export PDF disponible via ReportLab.
\end{itemize}
\end{frame}

\begin{frame}{Planification strategique}
\begin{itemize}
\item Calcul TRS:
\[
TRS = \frac{production\_reelle}{capacite\_theorique} \times 100
\]
\item Statuts:
\begin{itemize}
\item Critique si TRS < 80
\item Acceptable si 80 <= TRS < 90
\item Optimal si TRS >= 90
\end{itemize}
\item Generation d'une roadmap 3 ans dynamique.
\item La roadmap s'adapte a l'historique NLP (ex: motifs moteur/roulement).
\end{itemize}
\end{frame}

\begin{frame}{Orchestration des equipes}
\begin{itemize}
\item Donnees initiales: liste de techniciens, competences, disponibilite, charge.
\item Algorithme d'affectation:
\begin{enumerate}
\item filtrer techniciens disponibles avec la competence requise
\item trier par charge croissante
\item assigner le technicien le moins charge
\end{enumerate}
\item Si aucun profil compatible: statut Echec et mise en attente.
\item Visualisation en tableau de la charge de travail en temps reel.
\end{itemize}
\end{frame}

\begin{frame}{Persistence, historique et tracabilite}
\begin{itemize}
\item Base locale SQLite: table \texttt{events(id, event\_type, payload, created\_at)}.
\item Journalisation des actions utilisateur:
\begin{itemize}
\item chargement manuel
\item analyse NLP
\item generation prompt
\item export PDF
\item mise a jour planification
\item affectation taches
\end{itemize}
\item Onglet Historique: lecture des 50 derniers evenements.
\item Logs applicatifs: console + fichier rotatif (1 MB, 3 backups).
\end{itemize}
\end{frame}

\begin{frame}{Ce que le projet ne fait pas encore}
\begin{itemize}
\item Pas d'authentification utilisateur.
\item Pas de backend API operationnel expose.
\item Pas d'integration IoT active (MQTT/OPC-UA) dans la version actuelle.
\item Pas de modele de prediction ML entrainable en production.
\item Pas encore de packaging installeur finalise (MSI/Inno Setup).
\end{itemize}
\end{frame}

\begin{frame}{Stack technique reelle}
\begin{tabular}{ll}
\toprule
Couche & Technologies utilisees \\
\midrule
Desktop UI & Tkinter, ttk \\
Visualisation & Matplotlib, Seaborn, NumPy \\
NLP & spaCy, sentence-transformers \\
Documents & PyPDF2, python-docx \\
PDF & ReportLab \\
Donnees locales & SQLite (sqlite3) \\
Logs & logging + RotatingFileHandler \\
\bottomrule
\end{tabular}
\end{frame}

\begin{frame}{Conclusion}
\begin{itemize}
\item AURA est deja un socle desktop fonctionnel orientee maintenance industrielle.
\item Le code actuel privilegie l'usage local, la visualisation, l'analyse semantique et la traçabilite.
\item Les prochaines etapes les plus logiques sont: auth locale, theming UI avance, packaging executable Windows.
\end{itemize}
\end{frame}

\end{document}
